\documentclass[../main.tex]{subfiles}
\begin{document}

Chương này trình bày quá trình phân tích và thiết kế hệ thống ứng dụng hỗ trợ chăm sóc sức khỏe và dinh dưỡng tích hợp AI. Nội dung chương đi từ xác định yêu cầu chức năng theo 7 phân hệ nghiệp vụ và yêu cầu phi chức năng, phân tích tác nhân và ca sử dụng (Use Case), thiết kế kiến trúc tổng thể hệ thống, thiết kế cơ sở dữ liệu quan hệ trên PostgreSQL cho đến đặc tả API, giao tiếp thời gian thực và luồng giao diện người dùng.

\section{Xác định yêu cầu hệ thống}
\label{sec:3.1}

\subsection{Yêu cầu chức năng (Functional Requirements)}

Hệ thống được xây dựng như một trợ lý sức khỏe cá nhân toàn diện, tích hợp theo dõi dinh dưỡng, vận động, gợi ý thực đơn và tính năng tạo động lực. Bảng \ref{tab:ch3_fr} tổng hợp các yêu cầu chức năng chính theo từng phân hệ.

\begin{table}[H]
    \centering
    \caption{Danh sách yêu cầu chức năng của hệ thống}
    \label{tab:ch3_fr}
    \arrayrulecolor{primary}
    \rowcolors{2}{primary!10}{white}
    \begin{tabular}{|p{0.11\textwidth}|p{0.43\textwidth}|p{0.14\textwidth}|p{0.22\textwidth}|}
        \hline
        \rowcolor{primary!35}
        \textbf{Mã} & \textbf{Mô tả yêu cầu chức năng} & \textbf{Độ ưu tiên} & \textbf{Phân hệ} \\
        \hline
        FR-01 & Đăng ký, đăng nhập (Email / Google OAuth2) và khôi phục mật khẩu. & Cao & Tài khoản \& Bảo mật \\
        \hline
        FR-02 & Thiết lập thể trạng, mục tiêu cân nặng, thực phẩm dị ứng; tự động tính BMI, TDEE và calo mục tiêu hằng ngày. & Cao & Hồ sơ sức khỏe \\
        \hline
        FR-03 & Ghi nhật ký bữa ăn theo buổi (Sáng, Trưa, Tối, Bữa phụ) và tự tạo món ăn tùy chỉnh. & Cao & Dinh dưỡng \& Bữa ăn \\
        \hline
        FR-04 & Chụp/tải ảnh nhận diện món ăn Việt Nam bằng AI (FastAPI + ONNX), tự động tính calo/macro. & Cao & AI Recognition \\
        \hline
        FR-05 & Tra cứu dữ liệu món ăn và quét mã vạch thực phẩm đóng gói. & Cao & Dinh dưỡng \& Bữa ăn \\
        \hline
        FR-06 & Gợi ý thực đơn cá nhân hóa theo mục tiêu calo, dị ứng, ngân sách và Chatbot AI tư vấn theo lịch sử ăn uống. & Cao & Kế hoạch \& Trợ lý AI \\
        \hline
        FR-07 & Ghi nhận tập luyện, lịch sử cân nặng; đồng bộ bước chân, nhịp tim qua Google Fit / Apple Health. & Cao & Vận động \& Cơ thể \\
        \hline
        FR-08 & Dashboard thời gian thực; biểu đồ calo, macro, cân nặng (Ngày/Tuần/Tháng) và dự đoán xu hướng đạt mục tiêu. & Trung bình & Báo cáo \& Thống kê \\
        \hline
        FR-09 & Theo dõi streak hoàn thành mục tiêu, thử thách, huy hiệu và bảng xếp hạng bạn bè thời gian thực. & Trung bình & Streak \& Động lực \\
        \hline
        FR-10 & Nhắc nhở thông minh (uống nước, ghi log, cảnh báo vượt calo) qua FCM và hỗ trợ ngoại tuyến đồng bộ tự động. & Trung bình & Tiện ích \& Hệ thống \\
        \hline
    \end{tabular}
    \arrayrulecolor{black}
    \rowcolors{2}{}{}
\end{table}

\subsection{Yêu cầu phi chức năng (Non-functional Requirements)}

\begin{itemize}
    \item \textbf{Hiệu năng:} Thời gian phản hồi các API nghiệp vụ trên Spring Boot dưới $300\text{ ms}$; dịch vụ AI (FastAPI + ONNX Runtime) nhận diện món ăn và ước lượng calo dưới $1.5\text{ giây}$.
    \item \textbf{Bảo mật:} Xác thực không trạng thái bằng JWT kết hợp Google OAuth2; kiểm soát truy cập API theo vai trò qua Spring Security; truyền tải dữ liệu qua HTTPS.
    \item \textbf{Khả năng hoạt động ngoại tuyến (Offline-first):} Ứng dụng di động lưu tạm nhật ký bữa ăn, nước uống và vận động khi mất kết nối mạng, tự động đồng bộ lên PostgreSQL khi có mạng trở lại.
    \item \textbf{Thời gian thực \& khả năng mở rộng:} Cập nhật bảng xếp hạng bạn bè và phản hồi Chatbot qua WebSocket; tách biệt Backend chính và AI Service thành các container độc lập.
\end{itemize}

\section{Phân tích tác nhân và Use Case}
\label{sec:3.2}

Hệ thống gồm các tác nhân chính: \textbf{Người dùng (User)}, \textbf{Quản trị viên (Admin)}, \textbf{Dịch vụ AI (FastAPI + ONNX)} và các \textbf{Dịch vụ bên ngoài} (Google OAuth2, Google Fit / Apple Health, FCM).

\begin{figure}[H]
    \centering
    \resizebox{0.96\textwidth}{!}{%
    \begin{tikzpicture}[
        actor/.style={draw=primary, rounded corners, fill=primary!15, minimum width=2.3cm, minimum height=0.9cm, align=center, font=\small\bfseries},
        uc/.style={draw=primary, ellipse, fill=primary!6, minimum width=4.3cm, minimum height=0.8cm, align=center, font=\small},
        link/.style={->, thick, color=primary}
    ]
        \draw[draw=primary, thick, rounded corners=10pt, fill=primary!2] (-2.9, -4.8) rectangle (3.5, 4.8);
        \node[font=\small\bfseries, color=primary] at (0.3, 4.45) {Hệ thống Chăm sóc Sức khỏe \& Dinh dưỡng AI};

        \node[uc] (uc1) at (0.3, 3.4) {Tài khoản \& Hồ sơ (BMI/TDEE)};
        \node[uc] (uc2) at (0.3, 2.2) {Nhận diện món ăn \& Quét mã vạch};
        \node[uc] (uc3) at (0.3, 1.0) {Nhật ký bữa ăn \& Thực đơn AI};
        \node[uc] (uc4) at (0.3, -0.2) {Chatbot tư vấn dinh dưỡng};
        \node[uc] (uc5) at (0.3, -1.4) {Vận động \& Đồng bộ sức khỏe};
        \node[uc] (uc6) at (0.3, -2.6) {Thống kê, Streak \& Bảng xếp hạng};
        \node[uc] (uc7) at (0.3, -3.8) {Quản lý dữ liệu món ăn \& Hệ thống};

        \node[actor] (user) at (-5.6, 0.8) {Người dùng\\(User)};
        \node[actor] (admin) at (-5.6, -3.8) {Quản trị viên\\(Admin)};
        \node[actor] (ai) at (6.2, 1.6) {AI Service\\(FastAPI + ONNX)};
        \node[actor] (ext) at (6.2, -1.4) {Google Fit /\\Apple Health / FCM};

        \draw[link] (user) -- (uc1);
        \draw[link] (user) -- (uc2);
        \draw[link] (user) -- (uc3);
        \draw[link] (user) -- (uc4);
        \draw[link] (user) -- (uc5);
        \draw[link] (user) -- (uc6);
        \draw[link] (admin) -- (uc7);
        \draw[link] (uc2) -- (ai);
        \draw[link] (uc4) -- (ai);
        \draw[link] (uc5) -- (ext);
        \draw[link] (uc6) -- (ext);
    \end{tikzpicture}%
    }
    \caption{Sơ đồ Use Case tổng quát của hệ thống}
    \label{fig:ch3_usecase}
\end{figure}

\section{Thiết kế kiến trúc tổng thể hệ thống}
\label{sec:3.3}

Hệ thống áp dụng kiến trúc phân tách dịch vụ gồm ứng dụng di động đa nền tảng (React Native + Expo), máy chủ nghiệp vụ chính (Spring Boot), dịch vụ suy luận AI độc lập (Python FastAPI + ONNX Runtime), cơ sở dữ liệu quan hệ (PostgreSQL) và dịch vụ lưu trữ đối tượng đám mây (Cloud Object Storage).

\begin{figure}[H]
    \centering
    \resizebox{0.98\textwidth}{!}{%
    \begin{tikzpicture}[
        comp/.style={draw=primary, rounded corners, fill=primary!12, minimum width=3.3cm, minimum height=1.1cm, align=center, font=\small\bfseries},
        db/.style={draw=primary, rounded corners, fill=primary!6, minimum width=2.9cm, minimum height=0.95cm, align=center, font=\small},
        arrow/.style={<->, thick, color=primary}
    ]
        \node[comp] (client) at (0, 3.4) {Mobile App (TypeScript)\\React Native + Expo};
        \node[comp] (backend) at (0, 1.1) {Backend Chính\\Java Spring Boot + Security};
        \node[comp] (ai) at (5.3, 1.1) {AI Service (Python FastAPI)\\YOLO ONNX + Chatbot/Thực đơn};
        \node[db] (pg) at (-2.4, -1.3) {PostgreSQL\\(Dữ liệu quan hệ)};
        \node[db] (cos) at (2.4, -1.3) {Cloud Object Storage\\(Ảnh món ăn, Avatar)};
        \node[db] (ext) at (-5.5, 2.2) {Google Fit / Apple Health\\FCM + Expo Push};

        \draw[arrow] (client) -- node[right, font=\scriptsize] {REST / WebSocket (JWT)} (backend);
        \draw[arrow] (backend) -- node[above, font=\scriptsize] {HTTP / JSON} (ai);
        \draw[arrow] (backend) -- (pg);
        \draw[arrow] (backend) -- (cos);
        \draw[arrow] (client) -- (ext);
        \draw[arrow] (backend) |- (ext);
    \end{tikzpicture}%
    }
    \caption{Kiến trúc tổng thể hệ thống chăm sóc sức khỏe và dinh dưỡng tích hợp AI}
    \label{fig:ch3_architecture}
\end{figure}

Luồng xử lý chức năng nhận diện món ăn và ghi nhật ký bữa ăn diễn ra theo 5 bước:
\begin{enumerate}
    \item Người dùng chụp ảnh món ăn trên ứng dụng React Native/Expo; ứng dụng gửi ảnh kèm JWT token tới Backend Spring Boot.
    \item Backend Spring Boot xác thực token, lưu ảnh lên Cloud Object Storage để lấy URL và chuyển tiếp dữ liệu ảnh sang AI Service (FastAPI).
    \item AI Service chạy suy luận mô hình ONNX (huấn luyện trên VietFood68), tính toán khẩu phần, tra cứu bảng dinh dưỡng và trả về kết quả JSON gồm danh sách món ăn, độ tin cậy và calo/macro ước lượng.
    \item Backend trả kết quả về ứng dụng di động để người dùng kiểm tra, điều chỉnh khẩu phần (nếu cần) và xác nhận lưu vào bữa Sáng/Trưa/Tối/Phụ.
    \item Khi lưu nhật ký, Backend cập nhật dữ liệu vào PostgreSQL, tính lại tiến độ calo ngày, cập nhật chuỗi streak và phát thông báo cảnh báo qua FCM nếu vượt ngưỡng calo.
\end{enumerate}

\section{Thiết kế cơ sở dữ liệu}
\label{sec:3.4}

Cơ sở dữ liệu trên PostgreSQL được thiết kế chuẩn hóa nhằm phục vụ đồng thời quản lý dinh dưỡng, vận động, lập kế hoạch bữa ăn và tạo động lực:

\begin{table}[H]
    \centering
    \caption{Mô tả các bảng dữ liệu chính trong PostgreSQL}
    \label{tab:ch3_database}
    \arrayrulecolor{primary}
    \rowcolors{2}{primary!10}{white}
    \begin{tabular}{|p{0.23\textwidth}|p{0.35\textwidth}|p{0.32\textwidth}|}
        \hline
        \rowcolor{primary!35}
        \textbf{Tên bảng} & \textbf{Các trường dữ liệu chính} & \textbf{Vai trò} \\
        \hline
        \texttt{User} & \texttt{id, email, password\_hash, google\_id, full\_name, avatar\_url, role} & Lưu tài khoản, thông tin xác thực JWT / Google OAuth2. \\
        \hline
        \texttt{UserHealthProfile} & \texttt{id, user\_id, gender, age, height\_cm, weight\_kg, goal\_weight, allergies, bmi, tdee, target\_kcal} & Lưu thể trạng, dị ứng, mục tiêu và chỉ số BMI/TDEE tính toán. \\
        \hline
        \texttt{WeightLog}, \texttt{WaterLog} & \texttt{id, user\_id, weight\_kg, water\_ml, recorded\_at} & Theo dõi diễn biến cân nặng và lượng nước uống hằng ngày. \\
        \hline
        \texttt{Food}, \texttt{UserFoodLog} & \texttt{id, class\_label, barcode, name\_vi, calories, protein, carbs, fat, meal\_type, servings, image\_url} & Dữ liệu 68 món VietFood68, mã vạch, món tùy chỉnh và nhật ký bữa ăn (qua AI ảnh/mã vạch). \\
        \hline
        \texttt{MealPlan}, \texttt{MealPlanFood}, \texttt{DailySummary} & \texttt{id, user\_id, plan\_date, budget, food\_id, total\_kcal, total\_macro} & Lưu thực đơn cá nhân hóa và tổng hợp chỉ số dinh dưỡng, vận động theo ngày. \\
        \hline
        \texttt{ActivityLog}, \texttt{HealthSyncLog} & \texttt{id, user\_id, activity\_type, duration\_min, calories\_burned, steps, heart\_rate, source} & Lưu lịch sử tập luyện và dữ liệu đồng bộ từ Google Fit / Apple Health. \\
        \hline
        \texttt{Conversation}, \texttt{ChatMessage}, \texttt{Notification} & \texttt{id, user\_id, conversation\_id, sender, content, type, is\_read, created\_at} & Lưu hội thoại Chatbot AI và thông báo đẩy nhắc nhở thông minh (FCM). \\
        \hline
        \texttt{Streak}, \texttt{Badge}, \texttt{UserBadge}, \texttt{Challenge}, \texttt{UserChallenge}, \texttt{Friendship} & \texttt{id, user\_id, friend\_id, current\_streak, badge\_id, challenge\_id, progress, status} & Theo dõi chuỗi ngày mục tiêu, huy hiệu, thử thách và kết bạn bảng xếp hạng. \\
        \hline
    \end{tabular}
    \arrayrulecolor{black}
    \rowcolors{2}{}{}
\end{table}

\section{Thiết kế API và giao diện người dùng}
\label{sec:3.5}

Hệ thống cung cấp các nhóm RESTful API (Spring Boot \& FastAPI) và kênh WebSocket như mô tả trong Bảng \ref{tab:ch3_apis}:

\begin{table}[H]
    \centering
    \caption{Đặc tả các nhóm API và kênh truyền thông chính}
    \label{tab:ch3_apis}
    \arrayrulecolor{primary}
    \rowcolors{2}{primary!10}{white}
    \begin{tabular}{|p{0.14\textwidth}|p{0.34\textwidth}|p{0.42\textwidth}|}
        \hline
        \rowcolor{primary!35}
        \textbf{Giao thức} & \textbf{Endpoint / Kênh} & \textbf{Chức năng} \\
        \hline
        POST & \texttt{/api/v1/auth/login}, \texttt{/oauth2/google} & Đăng nhập Email/Mật khẩu hoặc Google OAuth2, cấp phát JWT. \\
        \hline
        GET / PUT & \texttt{/api/v1/profile/health} & Cập nhật thể trạng, dị ứng và tự động tính BMI, TDEE, calo mục tiêu. \\
        \hline
        POST & \texttt{/api/v1/ai/recognize-food} & Gửi ảnh tới FastAPI (ONNX) nhận diện món ăn Việt Nam và trả về calo/macro. \\
        \hline
        GET / POST & \texttt{/api/v1/meals}, \texttt{/api/v1/meals/sync} & Ghi nhật ký bữa ăn, tra cứu mã vạch và đồng bộ dữ liệu offline. \\
        \hline
        POST & \texttt{/api/v1/ai/meal-plan}, \texttt{/ai/chatbot} & Gợi ý thực đơn theo calo/dị ứng/ngân sách và hỏi đáp với Chatbot AI. \\
        \hline
        POST / GET & \texttt{/api/v1/workouts/sync}, \texttt{/analytics} & Đồng bộ Google Fit / Apple Health, thống kê và dự đoán xu hướng cân nặng. \\
        \hline
        WebSocket & \texttt{/ws/leaderboard}, \texttt{/ws/chat} & Cập nhật bảng xếp hạng bạn bè và phản hồi Chatbot thời gian thực. \\
        \hline
    \end{tabular}
    \arrayrulecolor{black}
    \rowcolors{2}{}{}
\end{table}

Về mặt giao diện người dùng (UI/UX) trên ứng dụng di động React Native, hệ thống được tổ chức thành 5 nhóm màn hình chính: (1) Màn hình Đăng nhập / Thiết lập hồ sơ sức khỏe; (2) Dashboard tổng quan theo thời gian thực (calo, macro, nước uống, bước chân, streak); (3) Màn hình Chụp ảnh nhận diện AI \& Quét mã vạch; (4) Màn hình Lập kế hoạch thực đơn \& Chatbot tư vấn AI; (5) Màn hình Thống kê biểu đồ, dự đoán cân nặng và Bảng xếp hạng bạn bè.

\end{document}
